%! TeX program = lualatex
% One page: the whole decision map for limits at infinity, in the style of
% mindmap-limits.tex -- three branches off the root, the same node shapes,
% colours and connector labels.  The box of known limits in the lesson
% (lim e^x, lim 1/x^p, e^x beats x^p, the polynomial sign rule) is deliberately
% NOT a branch here: this map is about what to do, not what to recall.
% `multi' is what makes each mmpage a page of its own -- the [tikz] option
% would do that for a bare tikzpicture, but it leaves the picture ungrouped
% and \resizebox below then fails with a missing \endgroup.
\documentclass[multi=mmpage]{standalone}
\input{../typography.tex.preamble}

% DO NOT use the tikz preamble: its `basic curve' styles pull in pgfplots, and
% enabling `trig format=rad' there makes this picture die with a dimension too
% large error.
% \input{../tikz.tex.preamble}
\usepackage{graphicx}   % for \resizebox
\usepackage{tikz}

\usetikzlibrary{mindmap, backgrounds, calc}

% The case that ends well: a finite answer, which is the one that names an
% asymptote.  WesternSpring on its own is a neon lime; muting it through grey
% keeps it in the palette without shouting.
\colorlet{finitecase}{WesternSpring!55!gray!45}
% Same hue at full strength, for the arrow joining two finitecase nodes.
\colorlet{finiteline}{WesternSpring!55!gray}
% The answer is zero -- still a finite answer, but worth its own fill so the
% three outcomes of a dominant-term analysis read as three.  WesternSky neat
% is a neon cyan, so mute it through grey the same way.
\colorlet{zerocase}{WesternSky!60!gray!45}
% The named theorem this map tells you to reach for.  Its own colour marks it
% as a different kind of thing: a theorem you invoke, not a case in a decision
% tree.  Same blue mindmap-limits.tex gives Squeeze and the IVT.
\colorlet{namedthm}{blue!20}

% Connector labels.  Neutral and italic rather than coloured: the node fills
% already carry the colour coding.
\colorlet{decision}{black!60}

\newenvironment{mmpage}{}{}

% Second line of a concept label.  Formulas stay in \displaystyle, as
% \everymath sets course-wide, so \lim subscripts sit underneath: the circles
% below are sized to allow for that.
\newcommand{\sub}[1]{\\[0.3ex]{\scriptsize\(#1\)}}
\newcommand{\note}[1]{\\[0.1ex]{\tiny #1}}

% The mindmap library gives a level-2 concept 1.5cm of text width and a
% level-3 one just 1cm, which hyphenates single words.  Widen them, and keep
% the manual line breaks short enough to fit what is set here.
\tikzset{
  mindmap base/.style={
    mindmap,
    grow cyclic,
    concept color=gray!20,
    every child node/.style={concept, minimum size=1.4cm, text width=1.7cm,
                             align=center, font=\scriptsize},
  },
  % Connector labels.  `mindmap' hands every node a large minimum size; that is
  % harmless when a label is centred on a point, but an edge-anchored one ends
  % up centimetres off its bar unless the box is shrunk back to its text.
  barlabel/.style={font=\footnotesize\itshape, text=decision, align=center,
                   inner sep=2pt, minimum size=0pt,
                   minimum width=0pt, minimum height=0pt},
}

% The level distances and sibling angles this map draws with.  Three branches
% rather than four, so level 1 fans by 120 degrees.
\tikzset{
  map levels/.style={
    mindmap base,
    level 1/.append style={level distance=5.5cm, sibling angle=120},
    level 2/.append style={level distance=4cm, sibling angle=60},
    level 3/.append style={level distance=3.4cm, sibling angle=44},
  },
}

% Two things about the angles on this page.
%
% (1) A `clockwise from' given to the root's children leaks down to every
%     deeper level, so without a correction every node's children fan in the
%     same absolute direction instead of pointing outward.  Each node
%     therefore restates its own fan, in the path right after it.
%     The value is (parent angle) - (children - 1) * (sibling angle) / 2.
%
% (2) Level distances and sibling angles differ per subtree, so they are set
%     in each child's option list, where they propagate down that subtree.

\begin{document}
% ===================================================================
%  The whole map -- limits at infinity
% ===================================================================
\begin{mmpage}
% The guided notes want a 7in-wide graphic.  The picture is drawn wider than
% that and scaled down to fit, so keep labels short: at this size every extra
% line of text costs legibility.
\resizebox{7in}{!}{%
\begin{tikzpicture}[map levels]
  % Branches at 150, 30 and -90: what the answer means, how to compute it, and
  % the moves that make the computation possible.
  % The axis in the root's label is the counterpart of the number line in
  % mindmap-limits.tex: no ticks, arrows running off both ways.
  \node [concept, minimum size=3.8cm, text width=3.2cm, font=\small] (root)
    {limit at infinity\\[0.8ex]
     % The nested picture inherits the mindmap's node sizing, which flings
     % these labels centimetres away; reset it here.
     \begin{tikzpicture}[baseline,
         every node/.style={inner sep=1.5pt, minimum size=0pt, minimum width=0pt,
                            minimum height=0pt, text width=, align=center}]
       \draw [<->, thin] (0,0) -- (2.6,0);
       \node [below] at (0.15,0) {\scriptsize\(-\infty\)};
       \node [below] at (2.45,0) {\scriptsize\(\infty\)};
     \end{tikzpicture}}
    [clockwise from=150]
    % ---------------------------------------------------------------
    % What the answer means.
    % ---------------------------------------------------------------
    child [concept color=gray!20,
           every node/.append style={text width=2.1cm},
           level 2/.append style={level distance=4.4cm, sibling angle=52}] {
      node (ha) [minimum size=2.6cm, text width=2.6cm, font=\footnotesize]
        {horizontal asymptotes\note{do \(+\infty\) and \(-\infty\)\\separately: at most two}}
      [counterclockwise from=98]
      child [concept color=finitecase] { node (finite) [text width=2.3cm]
        {a finite \(L\)\sub{y = L} } }
      child [concept color=pink!50] { node (infty) {\(+\infty\) or \(-\infty\)\note{no asymptote\\that side}} }
      child [concept color=supp!20] { node (dne) {no limit\note{oscillates: no asymptote}} }
    }
    % ---------------------------------------------------------------
    % How to compute it.
    % ---------------------------------------------------------------
    child [concept color=gray!20,
           level 2/.append style={level distance=4.2cm, sibling angle=56}] {
      % Both lines are broken by hand: left to itself TeX hyphenates
      % `dominant' rather than moving it down, and runs the note wider than
      % the blob it sits in.
      node (dominant) [minimum size=2.6cm, text width=2.6cm, font=\footnotesize]
        {analyze\\dominant terms\note{factor out the\\fastest-growing term}}
      [counterclockwise from=-26]
      child [concept color=pink!50] { node (toinfty) {\(\to \pm\infty\)\note{numerator wins}} }
      child [concept color=finitecase] { node (toconst) [text width=2.1cm]
        {\(\to\) a nonzero constant\note{same growth}} }
      child [concept color=zerocase] { node (tozero) {\(\to 0\)\note{denominator wins}} }
    }
    % ---------------------------------------------------------------
    % The moves that make the computation possible.
    % ---------------------------------------------------------------
    child [concept color=gray!20,
           level 2/.append style={level distance=4.2cm, sibling angle=46}] {
      node (moves) [minimum size=2.4cm, text width=2.2cm, font=\footnotesize]
        {before you factor}
      [counterclockwise from=-159]
      child { node (negx) [text width=2.4cm]
        {change of variable\sub{u = -x}\note{sends \(x \to -\infty\)\\to \(u \to \infty\)}} }
      child [concept color=supp!20] { node (expo) [text width=2.4cm]
        {factor out \(e^{x}\)\note{or \(u = e^{x}\);\\other bases: \(b^{x} = e^{x \ln(b)}\)}} }
      child [concept color=namedthm] { node (squeeze) {Squeeze\\Theorem\note{bounded times\\something small}} }
      child { node (graphit) [text width=2.3cm]
        {read the graph\note{trigs, inverse trigs,\\\(e^{x}\), \(\ln(x)\)}} }
    }
    ;

  % Labels riding the bars into a branch, rather than sitting in it as one
  % more node.  All of them are decisions or instructions, so they read as a
  % layer distinct from the concept names.
  \node [barlabel] at ($(root)!0.5!(dominant)$) {limit laws\\still apply};
  \node [barlabel] at ($(root)!0.5!(ha)$) {what the\\answer means};
  \node [barlabel] at ($(dominant)!0.5!(toconst)$) {compare\\growth rates};

  % The computation is what names the asymptote: the two branches are one
  % question asked twice.  Drawn on the background layer, in the finite case's
  % own colour, and shortened so the head clears the concept's blob outline.
  \begin{pgfonlayer}{background}
    % Bowed over the root rather than run straight across the top: a chord
    % from one branch's outermost child to the other's would cross every node
    % between them.  Head at the asymptote branch, because the value you
    % compute is what names the asymptote, not the other way round.
    \draw [concept connection, finiteline, ->, >=stealth, shorten >=4pt]
      (dominant) to [bend right=22] (ha);
    % Plain and grey, like the map's spine: these are the moves you make
    % before a dominant-term analysis, not a separate method.
    \draw [concept connection, gray!20, ->, >=stealth, shorten >=4pt]
      (moves) edge (dominant);
  \end{pgfonlayer}
\end{tikzpicture}
}
\end{mmpage}
\end{document}
